\documentclass[11pt,a4paper]{article}
\usepackage[margin=24mm]{geometry}
\usepackage{amsmath,amssymb}
\input{glyphtounicode}
\pdfgentounicode=1
\usepackage[hidelinks]{hyperref}
\hypersetup{hypertexnames=false}
\setlength{\parindent}{0pt}
\setlength{\parskip}{6pt}
\begin{document}
\begin{center}
{\Large C-P: conditional terminal-interface foundations}\\[4pt]
14 October 2026 --- analytical review, no selected replacement node
\end{center}

\textbf{Scope.} The user selected coupled point merging with local completion
at the choke. This note proves interface facts conditional on admissible
departure histories, not existence or physical correctness of a new service
law. The approved H1/T1 predecessor source/PDF are unchanged. No traffic
integration, parameter fit, numerical gradient or optimisation is run.

\section*{1. Retain H1 inputs and finite terminal histories}

Use the accepted road length $L=L_p+b$, initial footprint $[0,\ell]$,
total continuous mass $M$, fixed advice $v>0$, and initially empty terminal.
H1 produces conservative entrance histories $F_j$ for $j\in\{E,A,B\}$.
Let $w,k_j>0$, $H=k_jb$, $\tau=b/v$, $\tau_b=b/w$ and
$C(v)=vwk_j/(v+w)$. Histories are zero at negative times. The prospective
node must produce absolutely continuous, nondecreasing $Y_j$ satisfying
\begin{equation}
 P_j(t)=F_j(t-\tau),\qquad
 0\leq Y_j(t)\leq P_j(t),\qquad
 0\leq \dot Y_j(t)\leq C(v)
 \quad\hbox{a.e.}
 \tag{1}
\end{equation}
At an empty eligible queue the flow must also respect available arrivals;
this is included in the requirement $Y_j\leq P_j$. Define
\begin{equation}
 N_j=F_j-Y_j,\qquad q_j=P_j-Y_j\geq0.
 \tag{2}
\end{equation}
$N_j$ is physical stock; $q_j$ is eligible unserved mass. Unlike the old
nominal member, (1)--(2) do not require $Y_A=P_A$. They allow $q_A>0$.
Any extra service-memory variable is separate from vehicle mass. Traffic
that has crossed the objective boundary is not counted again as queue stock.

\section*{2. The conservative handoff bound does not require zero A delay}

Retain the sufficient domain
\begin{equation}
 M<H,\qquad 0\leq\dot F_j\leq C(v),\qquad F_j(t)\leq M.
 \tag{3}
\end{equation}
The accepted empty-link vacancy formula is
\begin{equation}
 V_j(t)=
 \begin{cases}
 k_jwt-F_j(t),&0<t\leq\tau_b,\\
 Y_j(t-\tau_b)+H-F_j(t),&t>\tau_b.
 \end{cases}
 \tag{4}
\end{equation}
For early times, $F_j(t)\leq C(v)t<k_jwt$ because $C(v)<k_jw$.
After $\tau_b$, nonnegative departures give
\begin{equation}
 V_j(t)\geq H-M>0.
 \tag{5}
\end{equation}
Thus the terminal entrance remains at its capacity under these conditional
first-order histories even if A or E departures temporarily stop. Inflow
within that capacity therefore does not restrict H1. This reuses the
predecessor's conservative proof, which never used $Y_A=P_A$.
It is not a general spillback solution or a proof that a microscopic
reference fits the chosen storage domain.

\newpage
\section*{3. Complete local TST includes every eligible lane queue}

Assume the eventual node clears the cohort in finite time $T$ and there are
no later arrivals or lost mass. The measured population remaining before
the choke is $M-\sum_jY_j(t)$. Therefore
\begin{align}
 \mathrm{TST}_{\rm choke}
 &=\int_0^T\left[M-\sum_jY_j(t)\right]dt\\
 &=\int_0^T\left[M-\sum_jP_j(t)\right]dt
   +\int_0^T\sum_jq_j(t)dt\\
 &=\frac{M(L-\ell/2)}{v}
   +\int_0^T(q_E+q_A+q_B)(t)dt.
 \tag{6}
\end{align}
The last equality uses unchanged H1's uniform initial footprint and common
ideal progress speed, its paired lateral conservation, and terminal
minimum travel $b/v$. Lateral exchange changes lane composition, not the
sum of minimum-travel completion histories. If B's old nonqueuing law is
retained, $q_B=0$ and the new term is A's queue area. Equation (6) is an
accounting identity, not a flow solver or a prediction of optimal advice.
Post-choke travel is excluded, while its service influence may remain in
the node state/history. Such influence is not added to (6) as fictitious mass.

\section*{4. Clearance needs a new service argument}

All H1 mass is minimum-travel eligible by $t_*=L/v$, so $\sum_jP_j(t_*)=M$.
Afterwards let $Q(t)=\sum_jq_j(t)$. A sufficient property of a future node is
\begin{equation}
 \sum_j\dot Y_j(t)\geq\mu_{\min}>0
 \quad\text{whenever }Q(t)>0\text{ and }t\geq t_*.
 \tag{7}
\end{equation}
Then $\dot Q=-\sum_j\dot Y_j\leq-\mu_{\min}$ while nonempty, and
\begin{equation}
 T\leq t_*+\frac{Q(t_*)}{\mu_{\min}}.
 \tag{8}
\end{equation}
No such lower bound has been established for a new coupled service law.
It is sufficient, not necessary: another valid law-specific clearance
proof can replace it. The old $R/\alpha$ guarantee cannot be copied when
the receiving/service state has changed. Likewise these interface facts
do not inherit a generic junction well-posedness theorem.

\section*{5. Conditional compact initial information}

Under continuing symmetry the admitted initial information reconstructs
\begin{equation}
 a_0=b_0=(\rho_0-e_0)/2.
 \tag{9}
\end{equation}
If all physical parameters and the new service initial state/history are
shared and known, H1 fixes the same $F_j$ for A+E and G. If the prospective
node is a unique autonomous causal map of those histories and its known
initial state, both representations give the same $Y_j$, queue histories
and (6). Predicted service memory does not by itself require an extra
initial observation. This is a conditional interface deduction; the new
map's existence/uniqueness and the old action/gradient certificate are
not proved here. An unknown initial memory would reopen G2.

\textbf{Source basis.}
\href{https://arxiv.org/pdf/1405.7080}{Jin, continuous LTM (2015),}
selected demand/vacancy and boundary-history sections support the physical
interface. All bounds, extended accounting and conditional reconstruction
above are deductions for the stated candidate inputs and admissible
departures, not source theorems about an unspecified coupled node.
\end{document}
